\documentclass[10pt,a4paper]{amsart}
\usepackage[margin=19mm]{geometry}
\usepackage{amsmath,amssymb,mathtools}
\usepackage[scr=rsfs,cal=euler]{mathalfa}
\usepackage{xfrac}
\usepackage[unicode,hidelinks,bookmarks=false]{hyperref}
\newcommand{\ie}{i.e.\ }
\newcommand{\cf}{cf.\ }
\newcommand{\resp}{respectively\ }
\newcommand{\Subsection}{\subsection}
\providecommand{\nobreakdash}{\nobreak\hbox{-}}
\setlength{\emergencystretch}{2em}
\multlinegap0pt
\usepackage{amssymb}
\usepackage{bm}
\usepackage{mathtools}
\usepackage{xcolor}
\usepackage{enumerate}
\newtheorem{lem}{Lemma}
\newtheorem{cor}{Corollary}
\newcommand{\An}{\mathrm{A}}
\newcommand{\Aut}{\mathrm{Aut}}
\newcommand{\Sn}{\mathrm{S}}
\def\ge{\geqs}
\newcommand{\geqs}{\geqslant}
\def\le{\leqs} 
\newcommand{\leqs}{\leqslant}
\newcommand{\minPart}{\Gamma_{k, |T|}}
\title{Greedy bases and relational complexity of diagonal type groups}
\author{Hong Yi Huang, Colva M. Roney-Dougal}
\date{Source: arXiv:2605.16032v1 (CC BY 4.0)}
\begin{document}
\maketitle

\noindent\emph{Source and licence.} Greedy bases and relational complexity of diagonal type groups. Version arXiv:2605.16032v1 (CC BY 4.0), distributed by arXiv under the Creative Commons Attribution 4.0 International licence, http://creativecommons.org/licenses/by/4.0/, as recorded in the arXiv licence metadata for this exact version and shown on the abstract page https://arxiv.org/abs/2605.16032v1. Full licence notice: \texttt{licenses/arxiv-2605.16032-CC-BY-4.0.txt}.\par\smallskip

\noindent\textit{Editorial scope.} Counted unit: the article's lem l:greedy\_two\_stab (source0.tex lines 1289--1292). The article's complete original proof of it is delivered with it (source0.tex lines 1294--1339, one \texttt{proof} environment of 46 lines). No mathematics is written, repaired or re-derived: the statement and the proof are the article's own lines, in the article's own notation, and no retained line is substituted or re-flowed.  Retained uncounted as the source-local context the proof needs: lines 1183--1188; lines 1241--1253; lines 1261--1263; lines 1274--1277.  Nothing was omitted from any retained span, so the package carries no omission record.  No retained line contains a citation, so the wrapper carries no bibliography and no bibliography entry.  No retained line contains a figure environment, an image-inclusion command or drawing code, so no image is delivered and the package declares no asset dependency.  Editorial formatting only, in addition: an AMS \texttt{amsart} preamble that restates the article's own declarations and macros used by the retained lines, namely the environments \texttt{lem}, \texttt{cor} on separate counters, together with the standard packages those lines need.  The retained environments print in this document as Lemma 1, Corollary 1 in order of appearance, whereas the article prints the same items with the numbering of its own full document.  The complete native source of the article as distributed in the arXiv e-print for this exact version is bundled unchanged as \texttt{sources/200702/source0.tex}.
\begin{lem}
	\label{l:part}
	Suppose that $k \ge |T| + 1$ and  let $\Pi$ be a partition of $[k]$ into $|T|$ %\leqs k-1$ 
    parts such that $\Pi \not\sim \minPart$,  $\Pi \not\sim \Sigma$ and $\Pi \not \sim [(m-1)^1, m^{|T|-2}, (m+1)^{1}]$.  Let $H$ be $\An_k$ or $\Sn_k$ as before. Then $|H_{(\Pi)}| > |H_{(\Sigma)}|$. Furthermore, $|H_{(\Sigma)}| \leqs 2|H_{(\minPart)}|$.
    %hh{[[$2$ to $3$ now, as $2$ is attained when $k = |T|+1$ or $|T|+2$ (with $m = 2$). This will be used in Lemma~\ref{l:greedy_two_stab} only, which I also make changes]]} hh{[[or we can say $|H_{(\Sigma)}| \leqs 2|H_{(\minPart)}|$, and replace all ``$\geqs 3|Q_{(\mathcal{P}^{\mathbf{t}})}| > $" by ``$> 2|Q_{(\mathcal{P}^{\mathbf{t}})}| \geqs $" in the proof of Lemma~\ref{l:greedy_two_stab}]]}
\end{lem}

\begin{cor}
	\label{c:stab_2-parts}
	Let $\mathbf{t} \in T^k$ and $\mathcal{P}^{\mathbf{t}}$ be as above. 
    \begin{itemize}\addtolength{\itemsep}{0.2\baselineskip}
        \item[{\rm (i)}] The group $Q_{(\mathcal{P}^{\mathbf{t}})}$ is a subgroup of the two-point stabiliser $G_{D , D\mathbf{t}}$.
        \item[{\rm (ii)}] If $G = W$ and $\mathcal{P}^{\mathbf{t}}$ has precisely two non-empty parts, labelled $1$ and $t$ and of distinct sizes, then 
	\begin{equation*}
	\begin{aligned}
	G_{D, D\mathbf{t}} &= \{(\varphi,\dots,\varphi)\sigma: \varphi\in C_{\Aut(T)}(t), \sigma\in P_{\mathcal{P}_1^{\mathbf{t}}, \mathcal{P}_t^{\mathbf{t}}}\} \cong C_{\Aut(T)}(t)\times P_{\mathcal{P}_1^{\mathbf{t}}}.
	\end{aligned}
	\end{equation*}
    \end{itemize}
\end{cor}

\begin{equation}\label{eq:wreath}
    \mbox{the group $T \wr \An_k$ is a subgroup of $G$, and so $\An_k \le Q \le \Sn_k$.}
\end{equation}

\begin{lem}
	\label{l:two_stab_sigma}
	Suppose $k\geqs |T|+1$. Then there exists an $\mathbf{s}\in T^k$ with $\mathcal{P}^{\mathbf{s}} = \Sigma$ such that $G_{D,D\mathbf{s}} = Q_{(\Sigma)}$.
\end{lem}

\begin{lem}
	\label{l:greedy_two_stab}
	Suppose $k\geqs |T|+1$, and let $\mathbf{t}\in T^k$ be such that $|G_{D,D\mathbf{t}}|$ is minimum. Then $\mathcal{P}^{\mathbf{t}}\sim\Sigma$ and $G_{D,D\mathbf{g}} = Q_{(\mathcal{P}^{\mathbf{t}})}$. 
\end{lem}

\begin{proof}
	By Lemma~\ref{l:two_stab_sigma},  there exists an $\mathbf{s}\in T^k$ with $\mathcal{P}^{\mathbf{s}} = \Sigma$ and $|G_{D,D\mathbf{s}}|  =  |Q_{(\Sigma)}|$. Let $\mathbf{t} = (g_1,\dots,g_k) \in T^k$ be such that $\mathcal{P}^{\mathbf{t}} \not \sim\Sigma$. It suffices to  show that $|G_{D,D\mathbf{t}}| > |Q_{(\Sigma)}|$.
    %whenever $\mathcal{P}^{\mathbf{g}}\not\sim\Sigma$. Write $\mathbf{g} = (g_1,\dots,g_k)\in T^k$.

    Corollary~\ref{c:stab_2-parts}(i) shows $G_{D,D\mathbf{t}} \geqs Q_{(\mathcal{P}^{\mathbf{t}})}$, and Lemma~\ref{l:part}  shows that if $\mathcal{P}^{\mathbf{t}}\not \sim \minPart$ and $\mathcal{P}^{\mathbf{t}} \not \sim [(m-1)^1, m^{|T|-2}, (m+1)^{1}]$ then 
	$|Q_{(\mathcal{P}^{\mathbf{t}})}| > |Q_{(\Sigma)}|$, which is equal to $|G_{D,D\mathbf{s}}|$.
	Thus, we may assume that either $\mathcal{P}^{\mathbf{t}}\sim \minPart$ (with five possibilities for $k$ modulo $|T|$), or $k = m|T|$ and $\mathcal{P}^{\mathbf{t}} \sim [(m-1)^1, m^{|T|-2}, (m+1)^{1}]$.
	%In each case it suffices to show there exists an element in $G_{D,D\mathbf{t}} \setminus Q_{(\mathcal{P}^{\mathbf{t}})}$, as then
	%\begin{equation*}
	%|G_{D,D\mathbf{g}}|\geqs 2|Q_{(\mathcal{P}^{\mathbf{g}})}| = 2|Q_{(\minPart)}| > |Q_{(\Sigma)}| = |G_{D,D\mathbf{s}}|
	%\end{equation*}
	%by Lemma~\ref{l:part} and Corollary~\ref{c:stab_2-parts}.
	
	%To prove that $G_{D,D\mathbf{g}} > Q_{(\mathcal{P}^{\mathbf{g}})}$, 
    First suppose  $k = m|T| +  \epsilon 1$, with $\epsilon = \pm$. Then $\mathcal{P}^{\mathbf{t}}\sim \minPart \sim [(m + \epsilon 1)^1,m^{|T|-1}]$. Without loss of generality, we may assume $|\mathcal{P}_1^{\mathbf{t}}| = m + \epsilon 1$. Then for any $x\in T$,
	\begin{equation*}
	\{t:|\mathcal{P}_t^{\mathbf{t}}| = m\} = T\setminus\{1\} = (T\setminus\{1\})^{x} = \{t:|\mathcal{P}_t^{\mathbf{t}}| = m\}^{x}.
	\end{equation*}
	Since $m \ge 2$ there exists $\sigma\in \An_k$ such that
	\begin{equation*}
	j^\sigma = i\mbox{ if and only if }g_j^{x} = g_i.
	\end{equation*}
	Hence \eqref{eq:wreath} shows that $(x,\dots,x)\sigma\in T \wr \An_k \leqs G$, and  we readily check that
	\begin{equation*}
	D^{(x, \ldots, x)\sigma} = D \quad \mbox{ and } \quad  D\mathbf{t}^{(x,\dots,x)\sigma} = D(g_{1^{\sigma^{-1}}}^{x},\dots,g_{k^{\sigma^{-1}}}^{x}) = D(g_1,\dots,g_k) = D\mathbf{t}.
	\end{equation*}
	Hence $|G_{D,D\mathbf{t}}| \geqs |T|\cdot |Q_{(\mathcal{P}^{\mathbf{t}})}| > 2|Q_{(\mathcal{P}^{\mathbf{t}})}| \geqs |Q_{(\Sigma)}|$ by Lemma~\ref{l:part}, as required. 
    %hh{[[For any $\tau\in Q_{(\mathcal{P}^{\mathbf{t}})}$, $(x,\dots,x)\sigma\tau$ gives another element in $G_{D,D\mathbf{t}}$, so the first equality follows from counting like this. We may omit explaining as it is just basic counting, if you are happy.]]}
	%$G_{D,D\mathbf{t}} >  Q_{(\mathcal{P}^{\mathbf{t}})}$, as required.  
    %The same argument can be applied to the case where $k = (m-1)|T|+1$, and we omit the details.
	
	Next, suppose $k = m|T| + \epsilon 2$, so $\mathcal{P}^{\mathbf{t}}\sim [(m + \epsilon 1 )^2,m^{|T|-2}]$, and assume that $|\mathcal{P}_1^{\mathbf{t}}| = m + \epsilon 1$. Let $z\in T\setminus\{1\}$ be the other element with $|\mathcal{P}_z^{\mathbf{t}}| = m + \epsilon 1$. Again, for any $x\in C_T(z)$,
	\begin{equation*}
	\{t:|\mathcal{P}_t^{\mathbf{t}}| = m\} = T\setminus\{1,z\} = (T\setminus\{1,z\})^{x} = \{t:|\mathcal{P}_t^{\mathbf{t}}| = m\}^{x}.
	\end{equation*}
	Hence again there exists $\sigma\in \An_k \le Q$ such that
	\begin{equation*}
	j^\sigma = i\mbox{ if and only if }g_j^{x} = g_i,
	\end{equation*}
	and it is easy to check that $(x,\dots,x)\sigma\in G_{D,D\mathbf{t}}$. Since no element of a non-abelian simple group has a  centraliser of order $2$,  
    %[[from an easy argument: the normaliser of a group of order $2$ in a Sylow $2$-subgroup has order strictly larger than $2$, as the order of a Sylow is larger than $2$ in a non-abelian simple group]], 
    it follows that $|G_{D,D\mathbf{t}}|\geqs |C_T(z)|\cdot |Q_{(\mathcal{P}^{\mathbf{t}})}| > 2|Q_{(\mathcal{P}^{\mathbf{t}})}| \geqs |Q_{(\Sigma)}|$ by Lemma~\ref{l:part}, as required.
	
	Finally, if $k = m|T|$, then either $\mathcal{P}^{\mathbf{t}}\sim \minPart \sim [m^{|T|}]$ or $\mathcal{P}^{\mathbf{t}}\sim [(m-1)^1,m^{|T|-2},(m+1)^1]$. We leave it as an exercise in each case to show that $|G_{D,D\mathbf{t}}| > 2|Q_{(\mathcal{P}^{\mathbf{t}})}| \geqs |Q_{(\Sigma)}|$.
	%find an element in $G_{D,D\mathbf{t}} \setminus Q_{(\mathcal{P}^{\mathbf{t}})}$.
\end{proof}

\end{document}
